\chapter{Testing and Verification}
\label{ch:testing}\index{testing}\index{verification}

A bare-metal system has nowhere to hide. There is no operating system to catch a
stray pointer, no test framework already on the board, and a bug in the field shows
up as a chip that silently locks up rather than a stack trace on a screen. So trust
here is built in layers, each catching what the one below it cannot. This chapter
gathers the testing story that the rest of the book uses in passing, from the
compiler that runs before anything else to the official Ada conformance suite
running on real silicon.

\section{Layer 0: the compiler and the type system}\index{type safety}

The first and cheapest test is the one you never run, because the compiler runs it
for you. A great deal of what a C program discovers at run time, an Ada program
rejects at compile time: a wrong-width register write, an out-of-range pin, a value
that violates a subtype predicate, a missing case alternative, a use of a resource
without holding it (Chapter~\ref{ch:packages}). This is not a substitute for testing,
but it moves a whole class of bug out of the test budget entirely, which is why the
chapters lean on types so hard.

\section{Layer 1: every listing compiles}\index{book!verify}

The code in this book is not pseudocode. The \texttt{book/verify} project is a
compile-only build whose sources are the worked listings from these chapters, built
against the \emph{real} HAL on the embedded profile for the Xtensa target. It has no
\texttt{Main}: \texttt{gprbuild} compiles every unit without linking, so a listing
that names a procedure that no longer exists, or passes the wrong type, fails the
build. The effect is that the prose cannot drift away from the code, because a
listing that stops compiling stops the book from building.

\section{Layer 2: host-native tests against an oracle}\index{host testing}

Much of the system is pure logic over an interface, and pure logic does not need the
board. The TLSF allocator (\texttt{tlsf\_core}, Chapter~\ref{ch:heap}), the ext4 block
logic (Chapter~\ref{ch:storage}), the Modbus framing (Chapter~\ref{ch:modbus}), the
routing table, and the language interpreters in the companion examples collection
(LISP, Lua, MicroPython, Forth) are all written target-free, so the identical Ada compiles and runs on
the host (x86) as an ordinary program. There it can be tested fast, in a debugger,
and against an \emph{oracle}: a second, independent implementation that is already
trusted.

The filesystem is the sharpest example. Its host tests build real images with the
Linux \texttt{mke2fs}, mutate them through the Ada code, and check every step against
the kernel's own tools: every read is byte-compared to the source file, every write
is checked with \texttt{e2fsck -fn}, and journal replay is cross-checked against the
kernel's recovery. The Modbus and routing tests drive the framing and the
longest-prefix logic against hand-built vectors. When a host test passes, the logic is right; only the
hardware wiring is then left to prove.

\section{Layer 3: known-answer tests for cryptography}\index{known-answer test}

Cryptography fails silently: a wrong result is not slower or noisier, it is just
wrong, and indistinguishable from right without a reference. So every primitive is
pinned with a \emph{known-answer test} (KAT), a fixed input whose correct output is
published. The examples \texttt{esp32s3\_aes\_gcm\_kat}, \texttt{esp32s3\_p256\_kat},
\texttt{esp32s3\_rsa\_kat}, \texttt{esp32s3\_sparknacl\_kat} and the X.509 set
(\texttt{x509\_kat}, \texttt{x509\_chain}, \texttt{x509\_policy},
\texttt{x509\_verify}) each carry vectors with a cited provenance (a NIST file and
count, or the exact script that generated them) so a reader can regenerate and trust
them. A KAT that still passes after a change is the proof that the change preserved
the maths.

\section{Layer 4: on-board self-tests}\index{self-test}\index{loopback}

A driver can be logically perfect and still mis-drive a pad. The cheapest hardware
test is one that needs no instrument and no second board, so the examples favour
\emph{reference-independent} self-tests. The TLV2556 ADC example reads its internal
zero, half, and full-scale references, which must come back as $0$, $2048$ and
$4095$ regardless of any external voltage. The loopback examples (\texttt{uart},
\texttt{i2c}, \texttt{i2s}, \texttt{rmt}, \texttt{twai}) wire a peripheral's output
back to its own input and check that what went out came back. The GDMA example copies
a buffer to itself and compares. When the BG95 modem refused to answer
(Chapter~\ref{ch:debug}), it was a UART single-pad loopback, output and input on the
same physical pin, that proved the transmit path was sound and put the fault on the
board, not the firmware. A self-test that passes proves the wiring, not just the
logic.

\section{Layer 5: conformance on real silicon}\index{ACATS}

The deepest layer tests the compiler and runtime themselves. The Ada Conformity
Assessment Test Suite (ACATS) is the official conformance suite for an Ada
implementation, and this project runs it \emph{on the chip}: each test is wrapped
into a bare-metal image, flashed, run, and graded from its serial output, with the
results matched structurally so a renamed test still scores. That sweep is the
subject of Chapter~\ref{ch:acats-sweep}; the point here is where it sits. The host
tests prove the application logic, the KATs pin the cryptography, the self-tests
prove the wiring, and the ACATS sweep proves that the language and the runtime the
whole stack stands on behave as the standard says, on the same silicon the product
ships on.

\section{Layer 6: proving the absence of errors with SPARK}\index{SPARK}\index{formal proof}

Every layer so far checks behaviour on inputs it was given. The top layer checks that a
whole category of failure cannot occur on \emph{any} input. Much of the pure logic ---
the ext4 serializers, the X.509, NMEA, and DNS parsers, the Modbus framing, the software
crypto and the conversion math --- is written in the SPARK subset of Ada and proved with
\texttt{gnatprove} to be free of run-time errors: no overflow, no buffer overrun, no
division by zero, no non-terminating loop, over the entire input space. For the code that
parses bytes chosen by a stranger, that is the difference between ``no fuzzer found a
crash'' and ``a crash is impossible.'' This is the subject of
Chapter~\ref{ch:spark}; it sits above the others because it rules out what they can only
sample.

\section{Putting the layers together}

No single layer is enough, and each is cheap where the next is dear. The compiler is
free and catches the most; the host tests are fast and oracle-checked but cannot see
a pad; the self-tests see the pad but need the board; the conformance sweep is slow
but proves the foundation; and the proofs, above them all, rule out an entire class of
run-time failure on inputs no test enumerates. A change to the filesystem runs the host
suite against \texttt{e2fsck} in seconds, a change to the curve re-runs a KAT, a change
to a driver re-runs its loopback on the board, a change to the runtime re-runs the
sweep, and a change to a parser re-runs its proof. Together they are why the code in
this book can be trusted to do on silicon what the prose says it does.
